% GENERATED FILE. Edit canonical lessons, then run npm run generate:books.
% canonical-source-hash: fnv1a64:2513831ac2c3ac7c
% canonical-lessons: SA-C491-jara, SA-C491-jiryati, SA-C491-mrtyuh, SA-C491-antyestih, SA-C491-sraddham

\chapter{Old Age and Death}
\label{ch:sa-old-age-and-death}
\hlchaptermodality{491}

\begin{chapteropening}
\textbf{By the end of this chapter:} \emph{I can talk about growing old and a death in the family.}

Complete the last lesson of chapter 491: I can talk about growing old and a death in the family.
\end{chapteropening}

\section[\texorpdfstring{\sk{जरा} (jarā)}{जरा (jarā)}]{\sk{जरा} (jarā) — old age}

\label{lesson:SA-C491-jara}



\begin{quote}
Before the new one: say the Sanskrit for a wrinkle, then the Sanskrit for appearance.
\end{quote}

\subsection*{You'll want to know: \sk{जरा}}
\textbf{\sk{जरा}} — \emph{jarā} — \textquotedblleft{}old age\textquotedblright{}.

\begin{quote}
\sk{जरायां} \sk{केशाः} \sk{श्वेताः} \sk{भवन्ति।}
\end{quote}

\begin{quote}
\emph{jarāyāṃ keśāḥ śvetāḥ bhavanti.}
\end{quote}

\begin{quote}
In old age the hair turns white.
\end{quote}

A feminine noun.

\begin{cousinweb}
\textbf{\sk{जरा}} (\emph{jarā}) comes from the root √jṝ, \textquotedblleft{}to grow old\textquotedblright{}; MW gives \textquotedblleft{}the act of becoming old, old age\textquotedblright{} (MW).
\end{cousinweb}

\begin{grammarlens}[title={What you've built}]
One more word for describing what you lived and felt.
\end{grammarlens}

\subsection*{Your turn}
\begin{itemize}
\raggedright
  \item \emph{Say it:} \emph{jarā}
  \item \emph{Say it:} \emph{jarā}, once more
  \item \emph{Say it:} \emph{jarāyāṃ keśāḥ śvetāḥ bhavanti.} — In old age the hair turns white.
  \item \emph{Recall:} say the Sanskrit for a wrinkle, then the Sanskrit for appearance, then say \emph{jarā} again
\end{itemize}

\begin{tcolorbox}[breakable,colback=teal!4,colframe=teal!35!black,title={Before you move on}]
What does \sk{जरा} mean? (\textquotedblleft{}Old age\textquotedblright{}.) Say it once more.
\end{tcolorbox}
\section[\texorpdfstring{\sk{जीर्यति} (jīryati)}{जीर्यति (jīryati)}]{\sk{जीर्यति} (jīryati) — wears out; grows old}

\label{lesson:SA-C491-jiryati}



\begin{quote}
Before the new one: say the Sanskrit for appearance, then the Sanskrit for old age.
\end{quote}

\subsection*{You'll want to know: \sk{जीर्यति}}
\textbf{\sk{जीर्यति}} — \emph{jīryati} — \textquotedblleft{}wears out; grows old\textquotedblright{}.

\begin{quote}
\sk{वस्त्रं} \sk{कालेन} \sk{जीर्यति।}
\end{quote}

\begin{quote}
\emph{vastraṃ kālena jīryati.}
\end{quote}

\begin{quote}
Cloth wears out with time.
\end{quote}

Past (it …): \sk{वस्त्रं} \sk{जीर्णवत्}, \emph{vastraṃ jīrṇavat}.

\begin{cousinweb}
\textbf{\sk{जीर्यति}} (\emph{jīryati}) comes from the root √jṝ, \textquotedblleft{}to grow old, to wear out\textquotedblright{}, whose ppp \emph{jīrṇa} is the A2 word \textbf{\sk{जीर्णः}} (MW).
\end{cousinweb}

\begin{grammarlens}[title={What you've built}]
One more word for describing what you lived and felt.
\end{grammarlens}

\subsection*{Your turn}
\begin{itemize}
\raggedright
  \item \emph{Say it:} \emph{jīryati}
  \item \emph{Say it:} \emph{jīryati}, once more
  \item \emph{Say it:} \emph{vastraṃ kālena jīryati.} — Cloth wears out with time.
  \item \emph{Recall:} say the Sanskrit for appearance, then the Sanskrit for old age, then say \emph{jīryati} again
\end{itemize}

\begin{tcolorbox}[breakable,colback=teal!4,colframe=teal!35!black,title={Before you move on}]
What does \sk{जीर्यति} mean? (\textquotedblleft{}Wears out; grows old\textquotedblright{}.) Say it once more.
\end{tcolorbox}
\section[\texorpdfstring{\sk{मृत्युः} (mṛtyuḥ)}{मृत्युः (mṛtyuḥ)}]{\sk{मृत्युः} (mṛtyuḥ) — death}

\label{lesson:SA-C491-mrtyuh}



\begin{quote}
Before the new one: say the Sanskrit for old age, then the Sanskrit for wears out.
\end{quote}

\subsection*{You'll want to know: \sk{मृत्युः}}
\textbf{\sk{मृत्युः}} — \emph{mṛtyuḥ} — \textquotedblleft{}death\textquotedblright{}.

\begin{quote}
\sk{पितामहस्य} \sk{मृत्युः} \sk{ह्यः} \sk{अभवत्।}
\end{quote}

\begin{quote}
\emph{pitāmahasya mṛtyuḥ hyaḥ abhavat.}
\end{quote}

\begin{quote}
Grandfather died yesterday.
\end{quote}

A masculine noun.

\begin{cousinweb}
\textbf{\sk{मृत्युः}} (\emph{mṛtyuḥ}) comes from the root √mṛ, \textquotedblleft{}to die\textquotedblright{}, the root of \textbf{\sk{म्रियते}}; MW gives \textquotedblleft{}death, dying\textquotedblright{} (MW).
\end{cousinweb}

\begin{grammarlens}[title={What you've built}]
One more word for describing what you lived and felt.
\end{grammarlens}

\subsection*{Your turn}
\begin{itemize}
\raggedright
  \item \emph{Say it:} \emph{mṛtyuḥ}
  \item \emph{Say it:} \emph{mṛtyuḥ}, once more
  \item \emph{Say it:} \emph{pitāmahasya mṛtyuḥ hyaḥ abhavat.} — Grandfather died yesterday.
  \item \emph{Recall:} say the Sanskrit for old age, then the Sanskrit for wears out, then say \emph{mṛtyuḥ} again
\end{itemize}

\begin{tcolorbox}[breakable,colback=teal!4,colframe=teal!35!black,title={Before you move on}]
What does \sk{मृत्युः} mean? (\textquotedblleft{}Death\textquotedblright{}.) Say it once more.
\end{tcolorbox}
\section[\texorpdfstring{\sk{अन्त्येष्टिः} (antyeṣṭiḥ)}{अन्त्येष्टिः (antyeṣṭiḥ)}]{\sk{अन्त्येष्टिः} (antyeṣṭiḥ) — funeral rites}

\label{lesson:SA-C491-antyestih}



\begin{quote}
Before the new one: say the Sanskrit for wears out, then the Sanskrit for death.
\end{quote}

\subsection*{You'll want to know: \sk{अन्त्येष्टिः}}
\textbf{\sk{अन्त्येष्टिः}} — \emph{antyeṣṭiḥ} — \textquotedblleft{}funeral rites\textquotedblright{}.

\begin{quote}
\sk{नद्याः} \sk{तीरे} \sk{अन्त्येष्टिः} \sk{भवति।}
\end{quote}

\begin{quote}
\emph{nadyāḥ tīre antyeṣṭiḥ bhavati.}
\end{quote}

\begin{quote}
The funeral rites take place on the riverbank.
\end{quote}

A feminine noun.

\begin{cousinweb}
\textbf{\sk{अन्त्येष्टिः}} (\emph{antyeṣṭiḥ}) is built from \emph{antya}, \textquotedblleft{}last\textquotedblright{}, and \emph{iṣṭi}, \textquotedblleft{}a sacrifice\textquotedblright{}, from the root √yaj; MW gives \textquotedblleft{}funeral sacrifice\textquotedblright{} (MW).
\end{cousinweb}

\begin{grammarlens}[title={What you've built}]
One more word for describing what you lived and felt.
\end{grammarlens}

\subsection*{Your turn}
\begin{itemize}
\raggedright
  \item \emph{Say it:} \emph{antyeṣṭiḥ}
  \item \emph{Say it:} \emph{antyeṣṭiḥ}, once more
  \item \emph{Say it:} \emph{nadyāḥ tīre antyeṣṭiḥ bhavati.} — The funeral rites take place on the riverbank.
  \item \emph{Recall:} say the Sanskrit for wears out, then the Sanskrit for death, then say \emph{antyeṣṭiḥ} again
\end{itemize}

\begin{tcolorbox}[breakable,colback=teal!4,colframe=teal!35!black,title={Before you move on}]
What does \sk{अन्त्येष्टिः} mean? (\textquotedblleft{}Funeral rites\textquotedblright{}.) Say it once more.
\end{tcolorbox}
\section[\texorpdfstring{\sk{श्राद्धम्} (śrāddham)}{श्राद्धम् (śrāddham)}]{\sk{श्राद्धम्} (śrāddham) — a rite for the ancestors}

\label{lesson:SA-C491-sraddham}



\begin{quote}
Before the new one: say the Sanskrit for death, then the Sanskrit for funeral rites.
\end{quote}

\subsection*{You'll want to know: \sk{श्राद्धम्}}
\textbf{\sk{श्राद्धम्}} — \emph{śrāddham} — \textquotedblleft{}a rite for the ancestors\textquotedblright{}.

\begin{quote}
\sk{पुत्रः} \sk{पितुः} \sk{श्राद्धं} \sk{करोति।}
\end{quote}

\begin{quote}
\emph{putraḥ pituḥ śrāddhaṃ karoti.}
\end{quote}

\begin{quote}
The son performs the rite for his late father.
\end{quote}

A neuter noun.

\begin{cousinweb}
\textbf{\sk{श्राद्धम्}} (\emph{śrāddham}) is built on \emph{śraddhā}, \textquotedblleft{}faith\textquotedblright{}; MW gives \textquotedblleft{}a ceremony in honour and for the benefit of dead relatives\textquotedblright{} (MW).
\end{cousinweb}

\begin{grammarlens}[title={What you've built}]
One more word for describing what you lived and felt.
\end{grammarlens}

\subsection*{Your turn}
\begin{itemize}
\raggedright
  \item \emph{Say it:} \emph{śrāddham}
  \item \emph{Say it:} \emph{śrāddham}, once more
  \item \emph{Say it:} \emph{putraḥ pituḥ śrāddhaṃ karoti.} — The son performs the rite for his late father.
  \item \emph{Recall:} say the Sanskrit for death, then the Sanskrit for funeral rites, then say \emph{śrāddham} again
\end{itemize}

\begin{tcolorbox}[breakable,colback=teal!4,colframe=teal!35!black,title={Before you move on}]
What does \sk{श्राद्धम्} mean? (\textquotedblleft{}A rite for the ancestors\textquotedblright{}.) Say it once more.
\end{tcolorbox}
